\documentclass[11pt,a4paper]{article}
\usepackage{turkos}

\begin{document}
\renewcommand{\progresslabel}{Stage V progress toward Turk-OS 2.0 (this phase in red)}
\phaseheader{16}{Portability: an Architecture Layer}{Find every place where x86 leaked into the generic kernel, hide the CPU and the board behind one interface, and prove that x86 still behaves exactly as before}{0}{20}

\ataglance{2--3 weeks}{3}{Phase 15 complete: Turk-OS 1.0 tagged, \code{make test} and the fuzzer pass}{The same Turk-OS on x86,
with all x86 code behind \code{include/arch.h}: a written leak audit in \code{docs/porting.md}, one trap frame with
accessors, \code{struct bootinfo} instead of Multiboot, the MMU behind \code{arch_mmu_*}, consoles and input that any
board can provide, \code{make ARCH=}, a 64-bit-clean generic core that compiles for \code{riscv32} and \code{aarch64},
a devicetree parser with host tests, and two more cross compilers ready for Phases 17 and 18.}

\section{Why this phase}

Turk-OS 1.0 is x86 from the first instruction to the last trap frame. Some of that is honest: \code{isr.s},
\code{gdt.c} and \code{paging.c} live in \code{kernel/arch/i386} because they can't be anything but x86. Much of it
leaked, though. The system call dispatcher reads \code{r->eax}, the page-fault handler reads CR2, \code{fork} writes
\code{tf->eax = 0}, the frame allocator walks the Multiboot memory map, \code{irq_restore} tests bit \code{0x200} of
EFLAGS, and \code{kmain} ends a test run by writing to port \code{0xF4}. Each of these lines works; it only shows up as
a leak on a second machine.

Stage V brings that second machine. Phase 17 ports Turk-OS to RISC-V (RV32 on QEMU's \code{virt} board, and optionally
your own Timur-RV32IMC core), and Phase 18 to 64-bit ARM (AArch64 on QEMU's \code{virt} board, and optionally a
Raspberry Pi~4). This phase writes no RISC-V or ARM code. It is a refactoring: it prepares the seams, and at its end
Turk-OS on x86 does exactly what version 1.0 did, with the same tests passing and benchmarks within a few per cent.

Porting sorts the kernel into three kinds of code. The \textbf{architecture} is the CPU: instruction set, privilege
levels, traps, MMU, atomics. The \textbf{board}, or platform, is everything around it: the devices, the firmware, and
how the kernel learns what memory and devices exist. Everything else is \textbf{generic}. \MOS\,\S12.3.1 draws the
same split as layers: hide the hardware first, then interrupts, context switching and the MMU, and above that the code
is mostly machine independent. The heart of this phase is the \textbf{leak audit} of Task~16.2, a written list of
every place where x86 shows through; every later task fixes entries on it. Stage V has its own progress scale, from 0
to 100\,\% toward Turk-OS 2.0; this phase takes you to 20\,\%.

\section{Concepts you need}

\subsection{Architecture, board and generic code}

Figure~\ref{fig:layers16} shows where Stage V ends up. Generic code calls the functions of \code{include/arch.h} and a
few inline helpers from per-architecture headers, and three implementations sit below that interface, each on its
board. On a PC the architecture and the board are entangled (port I/O is an x86 instruction, the PIC is part of every
PC), so Turk-OS keeps both in \code{kernel/arch/i386}. The ports separate them: their devices are found at boot, so
their drivers can live in \code{kernel/drivers} and serve any board that has the device.

\begin{figure}[htbp]
\centering
\begin{tikzpicture}[every node/.style={font=\sffamily\scriptsize},
  lab/.style={tknote, anchor=east, text width=1.15cm, align=right},
  arch/.style={tkbox, text width=4.25cm, minimum height=1.95cm},
  port/.style={arch, draw=tkViolet!85!black, fill=tkViolet!8},
  board/.style={tkbox, fill=white, text width=4.25cm, minimum height=1.55cm}]
  \node[tkok, text width=13.7cm, minimum height=1.0cm] (gen) at (0,0)
    {\textbf{generic kernel}: scheduler, VMM and copy-on-write, system calls, signals, VFS, TurkFS, buffer cache, tty,
     \code{kprintf}, ktests; drivers for devices that any board may have (16550 core, virtio-blk, RAM disk, block layer)};
  \node[tkhi, text width=13.7cm, minimum height=0.8cm] (ifc) at (0,-1.25)
    {\textbf{the interface}: \code{include/arch.h} (out-of-line functions), the inline headers \code{<arch/trapframe.h>},
     \code{<arch/irq.h>}, \code{<arch/memlayout.h>}, \code{<arch/mmu.h>}, \code{<arch/atomic.h>}, and \code{struct bootinfo}};
  \node[arch] (i386) at (-4.75,-2.95) {\textbf{kernel/arch/i386}\\GDT, IDT, \code{isr.s}, TSS, paging,\\\code{int 0x80}, PIC, PIT, VGA, ATA};
  \node[port] (rv) at (0,-2.95) {\textbf{kernel/arch/riscv32}\\\code{stvec}, Sv32, SBI calls, PLIC,\\\code{ecall}\quad(Phase 17)};
  \node[port] (a64) at (4.75,-2.95) {\textbf{kernel/arch/aarch64}\\\code{VBAR_EL1}, VMSAv8-64, GIC,\\generic timer, \code{svc}\quad(Phase 18)};
  \node[board] (pc) at (-4.75,-4.95) {\textbf{PC} (\code{qemu-system-i386})\\BIOS, GRUB, Multiboot;\\fixed I/O ports and addresses};
  \node[board] (rvb) at (0,-4.95) {\textbf{QEMU virt} (RISC-V)\\OpenSBI, devicetree; 16550, virtio\\later: Timur-RV32IMC};
  \node[board] (a64b) at (4.75,-4.95) {\textbf{QEMU virt} (Arm)\\devicetree; PL011, GICv2, virtio\\later: Raspberry Pi 4};
  \foreach \a in {i386,rv,a64} \draw[tkarrow] (\a.north) -- (\a.north |- ifc.south);
  \foreach \a/\b in {pc/i386,rvb/rv,a64b/a64} \draw[tkarrow] (\a) -- (\b);
  \node[lab] at (-6.95,0) {generic};
  \node[lab] at (-6.95,-1.25) {interface};
  \node[lab] at (-6.95,-2.95) {CPU};
  \node[lab] at (-6.95,-4.95) {board};
\end{tikzpicture}
\caption{The layers after Stage V. This phase builds the red interface and moves x86 below it; the violet backends
come in Phases 17 and 18.}
\label{fig:layers16}
\end{figure}

\subsection{Three machines, side by side}

Table~\ref{tab:three} is a map of what the interface hides; Phases 17 and 18 explain every RISC-V and ARM cell. What
stays the same matters as much: all three are little-endian, use 4\,KiB pages, have one unprivileged mode for programs,
and enter the kernel through one trap mechanism for system calls, faults and interrupts.

\begin{table}[htbp]
\centering\footnotesize\renewcommand{\arraystretch}{1.3}\setlength{\tabcolsep}{4pt}
\begin{tabularx}{\linewidth}{@{}>{\raggedright\arraybackslash}p{0.15\linewidth} >{\raggedright\arraybackslash}X >{\raggedright\arraybackslash}X >{\raggedright\arraybackslash}X@{}}
\toprule
\tkth{Mechanism} & \tkth{i386 (Turk-OS 1.0)} & \tkth{riscv32 (Phase 17)} & \tkth{aarch64 (Phase 18)} \\
\midrule
Firmware, entry & BIOS and GRUB; Multiboot hands over in protected mode, paging off & OpenSBI runs in M-mode and jumps to the kernel in S-mode, with the hart ID in \code{a0} and the devicetree's address in \code{a1} & QEMU's loader (or board firmware) starts the kernel at EL1, or at EL2 when virtualisation is enabled \\
Privilege levels & ring 0 and ring 3 & S-mode and U-mode, with M-mode (firmware) below & EL1 and EL0, with EL2 and EL3 above \\
Trap entry, fault address & IDT, one gate per vector; CR2 & \code{stvec}, one entry point (direct mode), cause in \code{scause}; \code{stval} & \code{VBAR_EL1}, 16 entries of 128 bytes, cause in \code{ESR_EL1}; \code{FAR_EL1} \\
Interrupt controller & two 8259 PICs & PLIC (claim and complete) & GIC (acknowledge and end of interrupt) \\
Timer & 8254 PIT & SBI timer call (TIME extension) and the \code{time} counter & Arm generic timer, frequency in \code{CNTFRQ_EL0} \\
Page tables & two levels, 32-bit entries, root in CR3 & Sv32: two levels, 32-bit entries, root in \code{satp} & VMSAv8-64, 4\,KiB granule, 39-bit addresses: three levels, roots in \code{TTBR0_EL1} and \code{TTBR1_EL1} \\
TLB maintenance & \code{invlpg}, or reload CR3 & \code{sfence.vma} & \code{tlbi} between \code{dsb} barriers, then \code{isb} \\
System call & \code{int 0x80}, number in EAX & \code{ecall}, number in \code{a7} & \texttt{svc \#0}, number in \code{x8} \\
Atomics & \code{xchg}, \code{lock} prefix & A extension: \code{amoswap.w}, or LR/SC & \code{ldaxr}/\code{stlxr} pairs (LSE atomics only from Armv8.1; the Cortex-A72 is Armv8.0) \\
Devices, discovery & port I/O and MMIO; Multiboot and fixed PC addresses & MMIO only; flattened devicetree & MMIO only; flattened devicetree \\
Word size & 32-bit, ILP32 & 32-bit, ILP32 & 64-bit, LP64 \\
\bottomrule
\end{tabularx}
\caption{The mechanisms behind the architecture interface, on the three targets of Stage V.}
\label{tab:three}
\end{table}

\subsection{One interface, three implementations}

There are two ways to keep one source tree for several machines. \MOS\,\S12.3.9 (\emph{Hiding the Hardware}) shows
\textbf{conditional compilation}: \code{config.h} defines \code{CPU} and \code{WORD_LENGTH}, and \texttt{\#if} blocks
select the code for each machine. Its lesson is orthogonality: test CPU-dependent things on the CPU macro and
word-length-dependent things on the word-length macro, so that a new CPU with a known word length needs no change to
the word-length code. A kernel full of \texttt{\#if (CPU == ...)} becomes hard to read, though, and every port touches
every file. Linux and MINIX 3 use \textbf{separate files behind one header} instead: a header declares what every
architecture provides, each architecture implements it in its own directory, and the build picks the directory.
Turk-OS does the same, and keeps \texttt{\#if} only for orthogonal features in Tanenbaum's sense:
\code{ARCH_HAS_ATOMICS} says whether atomic instructions exist, not which CPU this is.

\begin{conceptbox}{The architecture interface}
A fixed set of types, constants and functions that generic code may use and every architecture provides, in two
parts. \textbf{Inline headers} in \code{kernel/arch/$(ARCH)/include/arch/} hold what hot paths use or what needs the
architecture's types: the trap frame and its accessors, interrupt control, the memory layout, the PTE type. Being
\code{static inline}, they cost nothing at run time. \textbf{Out-of-line functions}, declared once in
\code{include/arch.h}, hold the rest: boot, the interrupt controller, the timer, context switching, page tables. A call
costs a few nanoseconds, nothing next to what these functions do. Task~16.12 checks both claims.
\end{conceptbox}

\subsection{The devicetree}

A PC kernel simply knows where things are: COM1 at port \code{0x3F8}, the disk at \code{0x1F0}, the PIC at
\code{0x20}, and GRUB reports the memory. Other boards have no such conventions, so their firmware passes a
\textbf{devicetree}: the machine as a tree of \textbf{nodes} with named \textbf{properties} (\SPEC Devicetree
Specification, chapter~2). Here is part of the tree that QEMU~10.2 generates for its RISC-V \code{virt} board:

\begin{lst}{plain}{QEMU riscv32 virt devicetree (excerpt in dtc's order; trimmed, comments added)}
/ {
    #address-cells = <0x02>;
    #size-cells = <0x02>;
    model = "riscv-virtio,qemu";
    memory@80000000 {
        device_type = "memory";
        reg = <0x00 0x80000000 0x00 0x8000000>;   // base 0x80000000, size 128 MiB
    };
    cpus { ... timebase-frequency = <0x989680>; ... };   // the time counter runs at 10 MHz
    chosen { stdout-path = "/soc/serial@10000000"; ... };
    soc {
        #address-cells = <0x02>;
        #size-cells = <0x02>;
        serial@10000000 {
            interrupts = <0x0a>;                   // source 10 of ...
            interrupt-parent = <0x03>;             // ... the node whose phandle is 3
            reg = <0x00 0x10000000 0x00 0x100>;
            compatible = "ns16550a";
        };
        plic@c000000 { phandle = <0x03>; compatible = "sifive,plic-1.0.0", "riscv,plic0"; ... };
    };
};
\end{lst}

\code{compatible} names the device, most specific first; a driver looks for a string it knows. \code{reg} lists
(address, size) pairs made of 32-bit cells, and the \emph{parent's} \texttt{\#address-cells} and \texttt{\#size-cells}
say how many cells each takes: two, 64 bits, even on this 32-bit board. \code{interrupts} is interpreted by the
controller that \code{interrupt-parent} names through its \code{phandle}, a numeric node ID. A few nodes have fixed
meanings (chapter~3): \code{/memory} lists RAM, \code{/reserved-memory} RAM the kernel must not touch, and
\code{/chosen} the command line (\code{bootargs}) and the console (\code{stdout-path}).

The firmware passes the tree as a \textbf{flattened devicetree} (FDT, or DTB), the binary format of chapter~5
(Figure~\ref{fig:fdt}). A 40-byte header with the magic number \code{0xd00dfeed} gives the offsets of three blocks.
The \textbf{structure block} is a stream of 32-bit tokens: \code{FDT_BEGIN_NODE} (1) and the node's name;
\code{FDT_PROP} (3), the value's length, the offset of the property's name in the \textbf{strings block}, and the
value; \code{FDT_END_NODE} (2); \code{FDT_NOP} (4); and \code{FDT_END} (9), all padded to four bytes. The
\textbf{memory reservation block} lists 64-bit (address, size) pairs to keep away from, ending with two zeros. Every
number is \textbf{big-endian}, and all three Turk-OS targets are little-endian, so every cell is byte-swapped as it is
read.

\begin{figure}[htbp]
\centering
\begin{tikzpicture}[blk/.style={cell, minimum height=0.9cm, text width=#1, fill=tkLight},
  tok/.style={cell, fill=white, minimum height=0.55cm, inner xsep=3pt}, every node/.style={font=\sffamily\scriptsize}]
  \node[blk=2.5cm, anchor=west, fill=tkRed!8] (h) at (0,0) {\textbf{header}, 40 bytes\\magic \code{0xd00dfeed}};
  \node[blk=3.3cm, anchor=west] (r) at (h.east) {\textbf{memory reservation}\\(address, size) pairs, then 0, 0};
  \node[blk=4.3cm, anchor=west, fill=tkBlue!8] (s) at (r.east) {\textbf{structure block}\\a stream of tokens (below)};
  \node[blk=3.0cm, anchor=west, fill=tkAmber!12] (t) at (s.east) {\textbf{strings block}\\property names, NUL-terminated};
  \node[tok, anchor=west] (k1) at (0,-1.45) {BEGIN\_NODE ""};
  \node[tok, anchor=west] (k2) at (k1.east) {PROP};
  \node[tok, anchor=west] (k3) at (k2.east) {\dots};
  \node[tok, anchor=west, fill=tkBlue!8] (k4) at (k3.east) {BEGIN\_NODE "memory@80000000"};
  \node[tok, anchor=west] (k5) at (k4.east) {PROP};
  \node[tok, anchor=west, fill=tkRed!8] (k6) at (k5.east) {PROP};
  \node[tok, anchor=west, fill=tkBlue!8] (k7) at (k6.east) {END\_NODE};
  \node[tok, anchor=west] (k8) at (k7.east) {\dots};
  \node[tok, anchor=west] (k9) at (k8.east) {END\_NODE};
  \node[tok, anchor=west] (k10) at (k9.east) {END};
  \draw[tkarrow, dashed] (s.south) -- (s.south |- k1.north);
  \node[tok, anchor=north west, fill=tkRed!8] (p1) at ([xshift=-1.6cm, yshift=-0.55cm]k6.south west) {3};
  \node[tok, anchor=west] (p2) at (p1.east) {len = 16};
  \node[tok, anchor=west] (p3) at (p2.east) {nameoff};
  \node[tok, anchor=west, text width=3.3cm] (p4) at (p3.east) {value: \code{00000000 80000000}\\\phantom{value: }\code{00000000 08000000}};
  \draw[tkarrow] (k6.south) -- (k6.south |- p1.north);
  \node[tknote, anchor=east] at (p1.west) {the \code{reg} property, one FDT\_PROP token:};
  \node[tknote, anchor=north east, align=right] at (p3.south east) {offset of "reg"\\in the strings block};
  \node[tknote, anchor=north west, align=left] at ([xshift=2pt]p4.south west) {base 0x80000000, size 128 MiB};
\end{tikzpicture}
\caption{The layout of a flattened devicetree, and the token sequence of the \code{reg} property of
\code{/memory@80000000}.}
\label{fig:fdt}
\end{figure}

\subsection{Word size: ILP32 and LP64}

i386 and riscv32 use the \textbf{ILP32} model: \code{int}, \code{long} and pointers are 32 bits. AArch64 uses
\textbf{LP64}: \code{int} stays at 32 bits, while \code{long}, pointers, \code{size_t} and \code{uintptr_t} grow to 64.
A pointer kept in a \code{uint32_t} or printed with \texttt{\%x} worked by accident on x86 and breaks on AArch64.
Table~\ref{tab:abi} lists the differences that bite; the last two even differ between the two 32-bit targets.

\begin{table}[htbp]
\centering\small\renewcommand{\arraystretch}{1.25}
\begin{tabularx}{\linewidth}{@{}>{\raggedright\arraybackslash}X c c c@{}}
\toprule
\tkth{Question} & \tkth{i386} & \tkth{riscv32} & \tkth{aarch64} \\
\midrule
\code{sizeof(long)}, \code{sizeof(void *)}, \code{sizeof(size_t)} & 4 & 4 & 8 \\
\code{sizeof(int)}; \code{sizeof(long long)} & 4; 8 & 4; 8 & 4; 8 \\
offset of \code{b} in \code{struct { uint32_t a; uint64_t b; }}, and the struct's size & 4, 12 & 8, 16 & 8, 16 \\
is plain \code{char} signed? & yes & no & no \\
\bottomrule
\end{tabularx}
\caption{The C data model on the three targets (checked with the compilers' output).}
\label{tab:abi}
\end{table}

What must \emph{not} change is what leaves the machine. TurkFS has only 8-, 16- and 32-bit fields, which all three
little-endian targets lay out the same way, so one disk image works everywhere; a \code{_Static_assert} on each
on-disk structure's size keeps it so. Elsewhere, use \code{uintptr_t} for an address in an integer, \code{size_t} for a
size, \code{uint64_t} for physical addresses and lengths from the firmware, and \code{uint8_t} for raw bytes. The best
auditor is the compiler: build the generic code for a target with another data model, and it reports every truncation
and format mismatch (Task~16.9).

\section{Reading plan}

\begin{readingplan}
\must & \MOS & \S12.3.9 \emph{Useful Techniques}: \emph{Hiding the Hardware}, \emph{Indirection} & 1005--1008 & Conditional compilation, orthogonal macros, indirection. \\
\must & \SPEC & Devicetree Specification (release v0.4), ch.~2 \emph{The Devicetree}: \S2.2 structure and conventions, \S2.3 standard properties, \S2.4 interrupts & -- & Nodes, \code{compatible}, \code{reg}, cells, \code{phandle}. \\
\must & \SPEC & Devicetree Specification, ch.~5 \emph{Flattened Devicetree (DTB) Format} & -- & The exact blob layout for the parser of Task~16.10. \\
\should & \SPEC & Devicetree Specification, ch.~3 \emph{Device Node Requirements}: \S3.4 \code{/memory}, \S3.5 \code{/reserved-memory}, \S3.6 \code{/chosen} & -- & What replaces the Multiboot information. \\
\should & \MOS & \S12.3.1 \emph{System Structure} (layered systems, Fig.~12-2) & 993--997 & Where the hardware-dependent layers belong. \\
\should & \OSDI & \S2.6.11 \emph{Hardware-Dependent Kernel Support} & 185--190 & How MINIX 3 kept x86-specific C code in a few files. \\
\should & \OSC & \S9.6 \emph{Example: Intel 32- and 64-bit Architectures}; \S9.7 \emph{Example: ARMv8 Architecture} & 379--384 & A preview of the MMUs of Stage V. \\
\could & \XVSIX & Ch.~1 \emph{Operating system interfaces}; ch.~2 \emph{Operating system organization} & -- & Meet xv6-riscv, Phase 17's reference kernel. \\
\end{readingplan}

For devicetree practice, the Linux kernel documentation page \emph{Linux and the Devicetree}
(\url{https://docs.kernel.org/devicetree/usage-model.html}) explains how a kernel uses the tree at boot.

\begin{minixbox}[Real-system lens: MINIX 3 on ARM, and Linux's arch/ tree]
The MINIX 3 of \OSDI (version 3.1) ran only on 32-bit x86, but \S2.6.11 explains that its hardware-dependent C code
was kept in a few files (\code{exception.c}, \code{i8259.c}, \code{protect.c}) to make porting easier. MINIX 3.3.0
(2014) was the first release with ARM support, for the BeagleBoard-xM and the BeagleBone (Cortex-A8, 32-bit ARMv7).
In its tree, \code{minix/kernel/arch/earm} stands next to \code{minix/kernel/arch/i386} with the same file names
(\code{head.S}, \code{mpx.S}, \code{klib.S}, \code{exception.c}, \code{protect.c}, \code{memory.c},
\code{arch_clock.c}, \code{kernel.lds}); \code{i8259.c} exists only on the x86 side, the TI board support only in
\code{earm/bsp/ti}. Linux does the same at scale: \code{arch/x86}, \code{arch/riscv} and \code{arch/arm64} each provide
headers under \code{include/asm}, and \code{include/asm-generic} supplies a default for whatever an architecture
doesn't override.
\end{minixbox}

\section{Build it}

\task{Two more cross compilers, two QEMU boards, dtc and a multi-architecture GDB}
Follow section~9 of \code{docs/setup.md}, \emph{Toolchains for the ports}. Section~9.2 installs the two emulators,
\code{dtc} and a GDB that knows all three architectures. Section~9.3 builds \code{riscv32-elf} and \code{aarch64-elf}
exactly like the \code{i686-elf} compiler of Phase~0, with the same binutils and GCC versions, the same prefix and one
build folder per target. Only GCC for RISC-V gets two more \code{configure} options,
\code{--with-arch=rv32imac_zicsr_zifencei} and \code{--with-abi=ilp32}, which make RV32IMAC and the \code{ilp32} ABI
its defaults. Run the checks of sections 9.5 and 9.6, then keep both boards' devicetrees for Task~16.10:

\begin{lst}{sh}{terminal}
mkdir -p tests/data
qemu-system-riscv32 -machine virt,dumpdtb=tests/data/virt-riscv32.dtb
qemu-system-aarch64 -machine virt,dumpdtb=tests/data/virt-aarch64.dtb -cpu cortex-a72
dtc -I dtb -O dts tests/data/virt-riscv32.dtb | less
\end{lst}
\verify{\code{riscv32-elf-gcc -v} shows \code{--with-arch=rv32imac_zicsr_zifencei --with-abi=ilp32} in its configure
line; \code{readelf -h} on a compiled object reports \code{ELF32} and \code{RISC-V} for one compiler, \code{ELF64} and
\code{AArch64} for the other. \code{qemu-system-riscv32 -machine virt -nographic -bios default} prints the OpenSBI
banner with \code{Domain0 Next Mode : S-mode}, and \code{dtc} decodes both blobs.}

\task{Find every place x86 leaked}
This is the heart of the phase. Before you move a single file, write down every place where code outside
\code{kernel/arch/i386} knows it runs on x86. Search for the usual suspects, then read every hit.

\begin{lst}{sh}{terminal}
grep -rnE --include='*.[ch]' -e 'struct regs|->(eax|ebx|ecx|edx|esi|edi|ebp|eip|cs|useresp)\b' \
  -e '\bcr[234]\b|invlpg|\b(inb|outb|inl|outl|insw|outsw)\(|\b(cli|sti|hlt)\b|pushf|0x200\b' \
  -e '0xC0000000|KERNEL_VMA|multiboot|rdtsc|pde_t|PG_[A-Z]|0xB8000|vga_|pic_|pit_|0xF4\b' \
  kernel lib include tests | grep -v '^kernel/arch/i386/'
grep -rnE --include='*.[ch]' '\(uint32_t\) *[&a-z_]|uint32_t +[a-z_]*(addr|ptr|len|size)\b' \
  kernel lib include | grep -v '^kernel/arch/i386/'
\end{lst}

Record each leak in \code{docs/porting.md}, a working document like the journal, with its kind: \textbf{arch} (a
register, an instruction, the page-table format), \textbf{board} (a device, a firmware structure, a fixed address) or
\textbf{word} size (a 32-bit type holding an address or a length), and the task that fixes it.

\begin{lst}{plain}{docs/porting.md (the first rows of a typical audit)}
| file:line                   | leak (kind)                          | fix (task)                 | done |
|-----------------------------|--------------------------------------|----------------------------|------|
| kernel/syscall/syscall.c:31 | arch: reads r->eax, r->ebx, ...      | tf_sysno, tf_arg (16.5)    | [ ]  |
| kernel/mm/pmm.c:40          | board: walks the Multiboot map       | struct bootinfo (16.6)     | [ ]  |
| kernel/syscall/uaccess.c:6  | word: uint32_t end = start + n       | uintptr_t (16.9)           | [ ]  |
\end{lst}

Some assumptions have no keyword for grep to find: the stack grows down, pages are 4\,KiB, the CPU is little-endian.
List them in a second table, \emph{assumptions kept}, and check that all three targets share each one (these three do).
\verify{Every hit of the two searches is in \code{porting.md}, with a kind and a fix, or marked there as a false
positive.}

\task{Move x86 into its corner}
Move everything that is x86 or PC into \code{kernel/arch/i386} with \code{git mv}, so the history follows each file:
the boot code and linker script (Phase~1's \code{boot/loader.s} and \code{link.ld} become \code{boot.s} and
\code{kernel.ld}), descriptor tables, interrupt stubs, PIC, PIT, TSS, paging, the context switch, and the drivers that
only exist on a PC (VGA text, PS/2 keyboard, ATA through ports). The x86 headers move from \code{include/arch/} to
\code{kernel/arch/i386/include/arch/}, and \code{include/} keeps only generic headers.

Generic code needs no change for this, thanks to the include path. Generic files keep writing
\texttt{\#include <arch/irq.h>}, and \code{CFLAGS += -Iinclude -Ikernel/arch/$(ARCH)/include} makes the compiler find
the header of whichever architecture is being built.
\verify{\code{include/arch/} no longer exists, \code{make clean && make test} passes, and
\code{grep -rn 'arch/i386' kernel lib include --include='*.c'} finds nothing outside \code{kernel/arch/}.}

\task{The architecture interface}
Write \code{include/arch.h}, the contract between the generic kernel and every port; Phases 17 and 18 implement
exactly these names. Most functions already exist on x86 under other names: rename them, and let generic code call
only these.

\begin{lst}{c}{include/arch.h}
#pragma once
#include <stdint.h>
#include <stdbool.h>
#include <bootinfo.h>
#include <arch/trapframe.h>    /* struct trapframe and the inline tf_ accessors */
#include <arch/mmu.h>          /* pte_t, paddr_t and struct arch_space          */

struct task;
struct context;                /* callee-saved registers, on the kernel stack */
struct addrspace;              /* include/vm.h                                 */
struct fault_info { uintptr_t addr; bool write, user, present, exec; };

/* boot, interrupts, time */
void     arch_early_init(struct bootinfo *bi);   /* console up, bootinfo filled */
void     arch_init(void);                        /* traps, interrupt controller  */
void     arch_irq_unmask(unsigned irq);
void     arch_irq_mask(unsigned irq);
void     arch_timer_init(unsigned hz);
uint64_t arch_cycles(void);                      /* TSC, time CSR, CNTVCT_EL0    */
uint64_t arch_entropy(void);                     /* seed for the stack canary    */
void     arch_fault_info(const struct trapframe *tf, struct fault_info *fi);

/* tasks, user mode, signals */
void     arch_task_init(struct task *t, void (*entry)(void *), void *arg);
void     arch_switch(struct context **save, struct context *load);
void     arch_user_entry(struct trapframe *tf, uintptr_t entry, uintptr_t sp);
_Noreturn void arch_return_to_user(struct trapframe *tf);
int      arch_signal_frame(struct trapframe *tf, int sig, uintptr_t handler, uintptr_t restorer);
int      arch_sigreturn(struct trapframe *tf);

/* address spaces; flags are VM_READ, VM_WRITE, VM_EXEC, VM_USER, VM_COW */
int      arch_mmu_init_space(struct addrspace *as);
void     arch_mmu_free_space(struct addrspace *as);
int      arch_mmu_map(struct addrspace *as, uintptr_t va, paddr_t pa, unsigned flags);
void     arch_mmu_unmap(struct addrspace *as, uintptr_t va);
bool     arch_mmu_lookup(struct addrspace *as, uintptr_t va, paddr_t *pa, unsigned *flags);
int      arch_mmu_protect(struct addrspace *as, uintptr_t va, unsigned flags);
void     arch_mmu_switch(struct addrspace *as);
void     arch_tlb_flush(uintptr_t va);

/* leaving */
_Noreturn void arch_test_exit(int failed);
_Noreturn void arch_halt(void);
\end{lst}

The inline headers complete the interface. \code{<arch/irq.h>} provides \code{unsigned long irq_save(void)},
\code{void irq_restore(unsigned long flags)}, \code{irq_enable()}, \code{irq_disable()} and \code{cpu_idle()}
(\code{sti; hlt} on x86); the flags type changes from Phase~9's \code{uint32_t} to \code{unsigned long}, the register
size everywhere. \code{<arch/atomic.h>} provides \code{atomic_xchg()} and \code{ARCH_HAS_ATOMICS} (1 on x86).
Phase~8's \code{switch_context(&prev->esp, next->esp)} becomes \code{arch_switch(&prev->ctx, next->ctx)} with a
\code{struct context *ctx}. The x86 assembly doesn't change: the saved ESP already points at the pushed EDI, ESI, EBX,
EBP and return address, which is all \code{struct context} is. \code{arch_task_init} builds a new task's first context
on its kernel stack: at the top, or just below \code{t->tf} when generic code has already put a trap frame there, as
\code{fork} does in Task~16.5. Keep the header small: a function joins it only when generic code really needs it.
\verify{\code{grep -rn 'switch_context\|qemu_exit\|enter_user_mode' kernel --include='*.c' | grep -v '^kernel/arch/'}
finds nothing, and \code{make test} passes.}

\task{One trap frame for generic code}
Rename Phase~3's \code{struct regs} to \code{struct trapframe}, with the same layout (\code{isr.s} builds it), and add
the accessors. From now on only code in \code{kernel/arch/i386} may name a field.

\begin{lst}{c}{kernel/arch/i386/include/arch/trapframe.h}
struct trapframe {                       /* was struct regs (Phase 3): same layout */
    uint32_t ds;
    uint32_t edi, esi, ebp, esp_at_pusha, ebx, edx, ecx, eax;
    uint32_t int_no, err_code;
    uint32_t eip, cs, eflags;
    uint32_t useresp, ss;                /* valid only for a trap from ring 3 */
};
_Static_assert(sizeof(struct trapframe) == 64, "isr.s depends on this layout");

static inline unsigned long tf_sysno(const struct trapframe *tf) { return tf->eax; }
static inline unsigned long tf_arg(const struct trapframe *tf, int n)
{
    const uint32_t r[6] = { tf->ebx, tf->ecx, tf->edx, tf->esi, tf->edi, tf->ebp };
    return r[n];                         /* the Linux i386 order */
}
static inline void tf_set_ret(struct trapframe *tf, long v) { tf->eax = (uint32_t)v; }
static inline bool tf_from_user(const struct trapframe *tf) { return (tf->cs & 3) == 3; }
/* tf_pc, tf_set_pc and tf_sp in the same style: eip, eip and useresp */
\end{lst}

Then rewrite the generic users. \code{syscall_dispatch(tf)}, which each port's trap code calls, sets
\code{current->tf}, checks \code{tf_sysno(tf)} against the table, calls the entry with six \code{unsigned long}
arguments, \code{tf_arg(tf, 0)} to \code{tf_arg(tf, 5)} (enough for every call on every port), and stores the result
with \code{tf_set_ret}. \code{fork} no longer needs Phase~11's \code{isr_return} label. It copies the parent's frame to
the top of the child's kernel stack, sets \code{child->tf} to it, calls \code{tf_set_ret(tf, 0)}, and then
\code{arch_task_init(child, fork_return, tf)}, where \code{fork_return(void *tf)} only calls
\code{arch_return_to_user(tf)}. The child starts as an ordinary kernel task whose first function returns to user mode.

The page-fault handler gets the same treatment: \code{arch_fault_info} fills a \code{struct fault_info} from CR2 and
the error code, and generic code decides what to do (Task~16.7). \code{exec} stays false on x86, because the error
code reports instruction fetches only with NX or SMEP, which Turk-OS doesn't use.
\verify{\code{grep -rnE -- '->(eax|ebx|eip|cs|useresp)' kernel lib --include='*.c' | grep -v '^kernel/arch/'} finds
nothing. The \code{fork_values} test and all seven evil programs of Task~10.9 behave as before.}

\task{Boot information without Multiboot}
Generic code must stop reading the Multiboot structure. Define what the kernel learns at boot, whatever the firmware,
and let each architecture fill it in. Addresses and lengths are \code{uint64_t} even on 32-bit targets: physical
memory can be larger than the address space, and the RISC-V devicetree uses 64-bit values too.

\begin{lst}{c}{include/bootinfo.h}
#define BOOTINFO_MAX_MEM  32
#define BOOTINFO_MAX_MODS 16

enum mem_type { MEM_USABLE = 1, MEM_RESERVED = 2 };

struct mem_region  { uint64_t base, length; enum mem_type type; };
struct boot_module { uint64_t start, end; char name[32]; };  /* physical; end is exclusive */

struct bootinfo {
    struct mem_region  mem[BOOTINFO_MAX_MEM];   unsigned nmem;
    struct boot_module mods[BOOTINFO_MAX_MODS]; unsigned nmods;   /* initrd and other modules */
    uint64_t           kernel_start, kernel_end;    /* physical range of the kernel image   */
    char               cmdline[256];
    const void        *fdt;                         /* the devicetree on the ports, or NULL */
    const char        *console;                     /* a hint: "com1", or /chosen stdout-path */
};
\end{lst}

On x86, \code{arch_early_init} calls \code{bootinfo_from_multiboot(bi, magic, mbi)} in
\code{kernel/arch/i386/multiboot.c}, which does what Phases~2, 5 and 10 did inline: it checks the magic
\code{0x2BADB002}, copies the memory map (flags bit~6) with type~1 as \code{MEM_USABLE} and everything else as
\code{MEM_RESERVED}, then the modules (bit~3) and the command line (bit~2), all read through \code{P2V}. It takes the
kernel's range from the linker symbols through \code{V2P}, sets \code{console} to \code{"com1"} and leaves \code{fdt}
\code{NULL}. The PMM becomes \code{pmm_init(const struct bootinfo *bi)}, and \code{kmain(struct bootinfo *bi)} never
sees Multiboot again.
\verify{The \code{mem} command prints the same regions and free-frame count as in version 1.0, with \code{-m 128} and
with \code{-m 512}. \code{grep -rln multiboot kernel lib include} lists only files under \code{kernel/arch/i386/}.}

\task{Memory layout and the MMU behind an interface}
Move the layout constants into \code{<arch/memlayout.h>}: \code{KERNEL_BASE} (\code{0xC0000000}, Phase~6's
\code{KERNEL_VMA} renamed), \code{USER_BASE} (\code{0x00400000}), \code{USER_STACK_TOP} (\code{KERNEL_BASE} on x86),
\code{PAGE_SIZE} (moved from \code{pmm.h}), \code{DIRECT_MAP_SIZE}, and \code{P2V} and \code{V2P} unchanged.

Then split the VMM. Generic code keeps the policy: \code{struct addrspace} (\code{include/vm.h}, holding the
\code{struct arch_space} of \code{<arch/mmu.h>}, on x86 the page directory), the ELF loader, \code{brk}, stack growth
and copy-on-write. Each architecture provides the mechanism, the \code{arch_mmu_*} functions, and translates the
generic flags into its PTE bits. On x86, \code{VM_WRITE} and \code{VM_USER} become \code{PG_WRITE} and
\code{PG_USER}, \code{VM_COW} becomes Phase~11's \code{0x200} (a bit left to the OS), and \code{VM_READ} and
\code{VM_EXEC} need no bit, because 32-bit paging has neither a read bit nor an NX bit. Every \code{arch_mmu_} call
leaves the TLB correct for the address it changed; \code{arch_mmu_init_space} copies the kernel half (PDEs
768--1023), and \code{arch_mmu_switch} replaces the CR3 load in \code{schedule()}. The copy-on-write case of
Phase~11 now reads as generic code, and Phases 17 and 18 inherit it unchanged:

\begin{lst}{c}{kernel/mm/fault.c (copy-on-write)}
static bool cow_fault(struct addrspace *as, uintptr_t va)    /* va: page-aligned */
{
    paddr_t old, copy;
    unsigned fl;
    if (!arch_mmu_lookup(as, va, &old, &fl) || !(fl & VM_COW))
        return false;
    fl = (fl & ~VM_COW) | VM_WRITE;
    if (pmm_refcount(old) == 1)                       /* last user: just take it back */
        return arch_mmu_protect(as, va, fl) == 0;
    if (!(copy = pmm_alloc_frame()))
        return false;
    memcpy(P2V(copy), P2V(old), PAGE_SIZE);
    arch_mmu_unmap(as, va);
    arch_mmu_map(as, va, copy, fl);
    pmm_unref(old);
    return true;
}
\end{lst}

\code{page_fault(tf)}, called by each port's trap code, gets a \code{struct fault_info} from \code{arch_fault_info},
tries \code{cow_fault} for a write to a present page below \code{KERNEL_BASE}, then Phase~11's lazy heap and stack
growth for a page that isn't present. Only then does it print the report (address, access, mode, \code{tf_pc(tf)}) and
send \code{SIGSEGV}, or panic for a fault in the kernel.
\verify{\code{grep -rnE 'pde_t|PG_[A-Z]|invlpg|cr3' kernel lib include | grep -v '^kernel/arch/i386/'} finds
nothing. The Phase~6 VMM tests, \code{cow_frames} and the \code{brk} test pass, and \code{crash null} still prints a
decoded report.}

\task{Consoles and input that any board can provide}
Phase~2's console layer called the VGA and serial drivers by name. Make it a registry, and give input one entry point:
the PS/2 keyboard's handler (x86 only) and a UART's receive interrupt (every board) both call \code{tty_input_char},
which feeds Phase~14's line discipline. Then split the serial driver, as the portability note of Task~2.6 suggested.
Everything about the 16550 (init, transmit, receive) moves to \code{kernel/drivers/uart16550.c} and reaches the
registers only through two accessors. The RISC-V \code{virt} board has the same chip, memory-mapped, so Phase~17 reuses
the file with a base address from the devicetree.

\begin{lst}{c}{include/console.h and kernel/drivers/uart16550.c (core)}
struct console {
    const char     *name;                        /* "vga", "com1", "uart0"          */
    void          (*putc)(struct console *c, char ch);
    struct console *next;
};
void console_register(struct console *c);        /* kprintf now also prints here    */
void tty_input_char(char c);                     /* every input device ends up here */

enum { UART_DATA = 0, UART_IER = 1, UART_FCR = 2, UART_LCR = 3, UART_MCR = 4, UART_LSR = 5 };

struct uart16550 {
    struct console con;
    uintptr_t      base;               /* an I/O port, or the registers' virtual address */
    unsigned       shift;              /* register n is at base + (n << shift)           */
    uint8_t      (*read)(const struct uart16550 *u, unsigned reg);
    void         (*write)(const struct uart16550 *u, unsigned reg, uint8_t val);
};

uint8_t uart_mmio_read(const struct uart16550 *u, unsigned reg)
{
    return *(volatile uint8_t *)(u->base + (reg << u->shift));
}

void uart16550_irq(struct uart16550 *u)          /* the port's IRQ handler calls this  */
{
    while (u->read(u, UART_LSR) & 0x01)          /* data ready: drain the receive FIFO */
        tty_input_char((char)u->read(u, UART_DATA));
}
\end{lst}

\code{uart_putc}, the \code{con.putc} of every 16550, waits for bit~5 of the line status register and writes the data
register, both through the accessors. The x86 side supplies port I/O accessors (\code{inb} and \code{outb} on
\code{base + reg}) in \code{kernel/arch/i386/uart_pio.c} and registers COM1 at \code{0x3F8}. Turn on serial
\emph{input} on x86 too, because it is the path the ports depend on: set bit~0 of the interrupt enable register
(received data available), write \code{0x0B} instead of Phase~2's \code{0x03} to the modem control register (OUT2
connects the UART's interrupt line to the PIC on a PC), unmask IRQ~4, and call \code{uart16550_irq} from its handler.
\verify{Boot output appears on the screen and in the \code{-serial stdio} terminal exactly as before. Keys typed into
the terminal reach \code{tsh} through IRQ~4, with editing and Ctrl+C working as on the PS/2 keyboard.}

\task{Make the generic code 64-bit clean}
Compile with \code{-Wall -Wextra -Wformat=2}, and mark \code{kprintf} with
\code{__attribute__((format(printf, 1, 2)))} so every format string is checked. GCC warns about casts between pointers
and integers of a different size by default (\code{-Wpointer-to-int-cast}, \code{-Wint-to-pointer-cast}). Teach
\code{kvprintf} the \code{l}, \code{ll} and \code{z} modifiers (\texttt{\%lx}, \texttt{\%llx}, \texttt{\%zu}). Fix what
the warnings show with the rules of the concepts section, and watch plain \code{char}: \code{if (c < 0)} is always
false on both ports (\code{-Wextra} says so, through \code{-Wtype-limits}), and a byte above \code{0x7F} used as an
index goes negative on x86.

On x86 alone, half of these mistakes stay silent, so add a target that compiles every generic file with both new
compilers and \code{-Werror}. Nothing is linked; the point is to let a 64-bit compiler read every line.

\begin{lst}{make}{Makefile (check-generic)}
CHECK_CFLAGS := -std=gnu11 -ffreestanding -O2 -Wall -Wextra -Werror -Wformat=2 \
                -Iinclude -Ikernel/arch/none/include
GENERIC_SRCS := $(filter-out kernel/arch/%,$(wildcard kernel/*.c kernel/*/*.c lib/*.c))

build/check/riscv32/%.o: %.c
	@mkdir -p $(dir $@)
	riscv32-elf-gcc -march=rv32imac_zicsr_zifencei -mabi=ilp32 $(CHECK_CFLAGS) -c $< -o $@
build/check/aarch64/%.o: %.c
	@mkdir -p $(dir $@)
	aarch64-elf-gcc -mgeneral-regs-only $(CHECK_CFLAGS) -c $< -o $@

check-generic: $(foreach t,riscv32 aarch64,$(patsubst %.c,build/check/$(t)/%.o,$(GENERIC_SRCS)))
	@echo "generic code compiles for riscv32 and aarch64"
.PHONY: check-generic
\end{lst}

\begin{decisionbox}{Which arch/ headers does check-generic use?}
The real RISC-V and ARM headers don't exist until Phases 17 and 18, so create \code{kernel/arch/none/include/arch/}:
compile-only stubs with plausible types (a trap frame of \code{unsigned long} registers, a 64-bit high-half
\code{KERNEL_BASE} when pointers are 64 bits, accessors that compile but mean nothing). Only \code{check-generic} uses
them, and nothing links against them. Writing the real port headers now would mean designing the RISC-V trap frame
before reading the privileged specification. Phase~17 switches the \code{riscv32} rule to its real headers; Phase~18
does the same for \code{aarch64} and deletes \code{arch/none}.
\end{decisionbox}
\verify{\code{make check-generic} ends with its message and no warnings. Put \code{uint32_t x = (uint32_t)ptr;} into
any generic file: the aarch64 rule now fails with GCC's \code{cast from pointer to integer of different size}. Take it
out again, and \code{make ARCH=i386 test} still passes.}

\task{Read a devicetree}
Write \code{lib/fdt.c}, a small read-only parser that both ports call from \code{arch_early_init} to fill
\code{struct bootinfo} and find their devices. Keep it free of kernel dependencies (only \code{<stdint.h>},
\code{<stddef.h>} and string functions that both \code{lib/string.c} and the host C library have), so it can be tested
on Linux as Phase~0 tested the warm-up library.

Its interface in \code{include/fdt.h} is \code{fdt_check(fdt)}, \code{fdt_next_node(fdt, node, &depth)},
\code{fdt_find_node(fdt, path)}, \code{fdt_getprop(fdt, node, name, &len)},
\code{fdt_next_compatible(fdt, after, compat)}, \code{fdt_reg(fdt, node, i, &base, &size)} and \code{fdt_be32(p)};
nodes are \code{int} offsets and $-1$ means ``none''. \code{fdt_check} compares the magic number, refuses a
\code{last_comp_version} above 17 (the version this parser knows), and checks that every block lies inside
\code{totalsize}. A node is the offset of its \code{FDT_BEGIN_NODE} token in the structure block, and everything else
is built on one walker that steps from node to node, keeps track of the depth, and never trusts a length:

\begin{lst}{c}{lib/fdt.c (the core)}
#define FDT_MAGIC 0xd00dfeedu
enum { FDT_BEGIN_NODE = 1, FDT_END_NODE = 2, FDT_PROP = 3, FDT_NOP = 4, FDT_END = 9 };

uint32_t fdt_be32(const void *p)                 /* byte by byte: any alignment, any CPU */
{
    const uint8_t *b = p;
    return (uint32_t)b[0] << 24 | (uint32_t)b[1] << 16 | (uint32_t)b[2] << 8 | b[3];
}
static uint32_t hdr(const void *fdt, int i) { return fdt_be32((const uint8_t *)fdt + 4 * i); }
static uint32_t align4(uint32_t n) { return (n + 3) & ~3u; }

int fdt_next_node(const void *fdt, int node, int *depth)
{
    const uint8_t *s = (const uint8_t *)fdt + hdr(fdt, 2);   /* off_dt_struct  */
    uint32_t end = hdr(fdt, 9), off, len;                    /* size_dt_struct */
    if (node < 0) {                                          /* the root node  */
        *depth = 0;
        return fdt_be32(s) == FDT_BEGIN_NODE ? 0 : -1;
    }
    off = node + 4 + align4(strlen((const char *)s + node + 4) + 1);   /* token, name */
    while (off + 4 <= end) {
        switch (fdt_be32(s + off)) {
        case FDT_BEGIN_NODE: (*depth)++; return (int)off;
        case FDT_END_NODE:   (*depth)--; off += 4; break;
        case FDT_NOP:        off += 4; break;
        case FDT_PROP:                                       /* len, nameoff, value */
            if (end - off < 12 || (len = fdt_be32(s + off + 4)) > end - off - 12)
                return -1;
            off += 12 + align4(len);
            break;
        default:             return -1;                      /* FDT_END, or corrupt */
        }
    }
    return -1;
}
\end{lst}

\code{fdt_getprop} walks the properties after a node's name up to its first child. \code{fdt_find_node} matches a
path one component per depth. \code{fdt_reg} reads the \emph{parent's} \texttt{\#address-cells} and
\texttt{\#size-cells} (the defaults are 2 and 1) and joins one or two cells into each 64-bit value. Test it in
\code{tests/host/test_fdt.c}, compiled with the host \code{gcc} and \code{-fsanitize=address,undefined} as in Phase~0,
against facts from the \code{dtc} output of the riscv32 blob: \code{/memory@80000000} has base \code{0x80000000} and
size 128\,MiB (QEMU's default \code{-m}); the first \code{ns16550a} node has base \code{0x10000000} and
\code{interrupts} 10; \code{/cpus} has a \code{timebase-frequency} of 10\,000\,000; there are eight
\code{virtio,mmio} nodes; and \code{/chosen} has a \code{stdout-path}. Write the same kind of test for the aarch64 blob
(its memory node, the PL011, the timer). Finally, try to break the parser: flip a byte of the magic, make
\code{off_dt_struct} plus \code{size_dt_struct} point past \code{totalsize}, set a property's length to
\code{0xFFFFFFF0}. Each must give an error, never a crash or an endless loop.
\verify{The host tests pass on both blobs under the address and undefined-behaviour sanitizers, the corrupted blobs
are rejected, and \code{lib/fdt.c} also compiles in \code{make check-generic}.}

\task{One Makefile, three architectures}
Make the architecture a build variable. What differs per architecture moves into \code{kernel/arch/$(ARCH)/arch.mk}:
cross prefix, extra flags, sources, linker script, debugger, how to run QEMU, and how a test run reports success. The
\code{GDB} line from \code{docs/setup.md}, section~4, moves there too (section~9.4). Output goes to
\code{build/$(ARCH)/}, so objects of different architectures never mix.

\begin{lst}{make}{Makefile (the parts that change), then kernel/arch/i386/arch.mk}
ARCH     ?= i386
ARCH_DIR := kernel/arch/$(ARCH)
BUILD    := build/$(ARCH)
ifeq ($(wildcard $(ARCH_DIR)/arch.mk),)
$(error unknown ARCH '$(ARCH)': there is no $(ARCH_DIR)/arch.mk)
endif
include $(ARCH_DIR)/arch.mk
CC     := $(CROSS)gcc
CFLAGS += $(ARCH_CFLAGS) -Iinclude -I$(ARCH_DIR)/include
OBJS   := $(addprefix $(BUILD)/,$(addsuffix .o,$(basename $(GENERIC_SRCS) $(ARCH_SRCS))))
# run, debug and test use $(QEMU) $(QEMU_FLAGS); test adds $(QEMU_TEST_FLAGS)
# and compares QEMU's exit status with $(TEST_PASS_STATUS)

# ---- kernel/arch/i386/arch.mk ----
CROSS            := i686-elf-
ARCH_SRCS        := $(wildcard kernel/arch/i386/*.c kernel/arch/i386/*.s)
LDSCRIPT         := kernel/arch/i386/kernel.ld
GDB              := $(call find_tool,i386-elf-gdb gdb)
QEMU             := qemu-system-i386
QEMU_FLAGS       := -kernel $(BUILD)/kernel.elf -serial stdio
QEMU_TEST_FLAGS  := -append "test" -display none -no-reboot \
                    -device isa-debug-exit,iobase=0xf4,iosize=0x04
# isa-debug-exit: QEMU exits with 2 * value + 1, and a passing run writes 0
TEST_PASS_STATUS := 1
# ... plus the NASM pattern rule for the .s files, which only x86 has
\end{lst}

Do the same for user space: \code{crt0.s} and the system call stubs move to \code{user/arch/i386/}, and user programs
get their own flags without \code{-Ikernel}. Phases 17 and 18 add the \code{arch.mk} files for \code{riscv32-elf-}
and \code{aarch64-elf-}; until then \code{make ARCH=riscv32} stops with the error.
\verify{\code{make} and \code{make ARCH=i386} build the same \code{build/i386/kernel.elf}; \code{make ARCH=i386 run},
\code{debug} and \code{test} work; \code{make ARCH=sparc} prints the error; nothing is written to \code{build/} outside
\code{build/i386/} and \code{build/check/}.}

\task{Prove that nothing broke}
A refactoring is finished when the behaviour is provably the same. Run Phase~15's safety net against the new tree
and compare with your v1.0 benchmark table.

\begin{enumerate}
  \item \code{make clean && make ARCH=i386 test}: every ktest and user-space test passes.
  \item Four fuzzers and the stress script of Task~15.1 for ten minutes: no panic, no growth in memory use.
  \item The benchmarks of Task~15.5, above all the null system call (now through \code{syscall_dispatch} and six
        \code{tf_arg} calls) and the context switch (now \code{arch_switch}).
  \item \code{make check-generic} and the FDT host tests, from now on before every commit.
\end{enumerate}

If the null system call got slower, read the machine code before guessing:
\code{i686-elf-objdump -d build/i386/kernel/syscall/syscall.o} must show no \code{call} to a \code{tf_} function (an
accessor declared without \code{static inline} ends up out of line). Record the numbers in \code{docs/porting.md}.
\verify{All four items pass, and every benchmark is within a few per cent of its Phase~15 value (repeat the runs and
compare medians before deciding that a difference is real).}

\section{Testing and debugging}

\begin{pitfalls}
A change to an x86 header has no effect, or two definitions clash & A stale \code{include/arch/} directory, or a leftover \code{-Iinclude/arch}, is found before \code{kernel/arch/$(ARCH)/include} & Delete the old directory; print the search path with \code{i686-elf-gcc -E -v} and check its order. \\
The child of \code{fork} crashes or returns garbage & \code{arch_task_init} built the first context over the trap frame at the top of the kernel stack, or code still writes \code{tf->eax} & Build the context below \code{t->tf} when it is set; grep for field names outside \code{kernel/arch}. \\
\code{fdt_check} fails; the magic reads \code{0xedfe0dd0} & The header was read as a little-endian \code{uint32_t} & Read every cell through \code{fdt_be32}. \\
RAM above 4\,GiB, or a large region, is missing or wrong on one target & A \code{uint64_t} base or length passed through a \code{uint32_t} or \code{uintptr_t} on the way into \code{bootinfo} & Keep firmware values 64-bit until the PMM clips them to what it can map. \\
Benchmarks several per cent slower & Accessors became function calls, or \code{arch_switch} gained work & Look for \code{call} in \code{objdump -d}; keep hot helpers \code{static inline}. \\
\end{pitfalls}

\section{Milestone}

\begin{milestonebox}[Portable core]
\begin{checklist}
  \item Generic code is free of x86: every row of the audit in \code{docs/porting.md} is done, and the searches of
        Task~16.2 find nothing outside \code{kernel/arch/i386}.
  \item \code{make ARCH=i386 test} passes; ten minutes of fuzzing cause no panic; the benchmarks are within a few
        per cent of Phase~15.
  \item \code{make check-generic} compiles every generic file for \code{riscv32} and \code{aarch64} without warnings.
  \item The FDT parser passes its host tests on both QEMU blobs under the sanitizers.
  \item \code{include/arch.h}, the trap frame accessors, \code{struct bootinfo}, \code{arch_mmu_*},
        \code{console_register} and \code{tty_input_char} are in place; serial input works on x86.
  \item Both new cross compilers, both QEMU boards, \code{dtc} and a GDB for all three architectures are installed and
        smoke-tested.
\end{checklist}
\end{milestonebox}

\begin{questionbox}
\begin{enumerate}
  \item What is the difference between the architecture and the board? Why does x86 keep both in one directory?
  \item Why are some parts of the interface inline headers and others functions in \code{include/arch.h}? What would
        happen to the null system call if \code{tf_arg} were a function?
  \item Where in Turk-OS was an \texttt{\#ifdef} most tempting, and what did you do instead?
  \item The devicetree is big-endian, and all three targets are little-endian. What would go wrong without
        \code{fdt_be32}, and how would you notice?
  \item Why can one TurkFS image work on all three targets, and which kinds of change to its structures would break
        that?
  \item Which subsystems needed no change at all, and what does that say about their design?
\end{enumerate}
\end{questionbox}

\section{Going further}

Add a fourth architecture that is not a machine: \code{kernel/arch/host} runs the generic kernel as a Linux process,
with \code{mmap} standing in for page tables, so the ktests of the heap, the scheduler and TurkFS run in a second under
the sanitizers (User-Mode Linux does this for all of Linux). Run \code{make check-generic} and the host tests from a
Git pre-commit hook. Generate the system call table, the user stubs of every architecture and \code{docs/syscalls.md}
from one list. Or add Clang to \code{check-generic}: \code{clang --target=riscv32-unknown-elf} and
\code{--target=aarch64-unknown-elf} need no toolchain build, and a second compiler finds different mistakes.

\nextphase{17}{The RISC-V Port}
\booklegend
\portlegend
\end{document}
